% BHU PYQ -- Manually transcribed & error-corrected
% Paper: BSC-06A : Ancillary Mathematics
% Exam: B.Sc. (Hons.) Semester II Examination, 2013-14 (Ancillary, Old Course)

\documentclass[11pt,a4paper]{article}
\usepackage[a4paper,top=2.5cm,bottom=2.5cm,left=2.8cm,right=2.8cm,headheight=14pt]{geometry}
\usepackage{amsmath,amssymb}
\usepackage[T1]{fontenc}
\usepackage[utf8]{inputenc}
\usepackage{lmodern}
\usepackage{microtype}
\usepackage{fancyhdr}
\usepackage{enumitem}
\usepackage{hyperref}
\usepackage{lastpage}
\usepackage{parskip}

\hypersetup{colorlinks=true,urlcolor=black,linkcolor=black,
  pdftitle={BSC-06A Ancillary Mathematics -- BHU B.Sc. Sem II 2013-14}}
\pagestyle{fancy}\fancyhf{}
\renewcommand{\headrulewidth}{0.4pt}\renewcommand{\footrulewidth}{0.4pt}
\fancyhead[L]{\small   \textit{Banaras Hindu University}}
\fancyhead[R]{\small   \textit{BSC-06A: Ancillary Mathematics}}
\fancyfoot[C]{\small Page  \thepage\ of \pageref{LastPage}}


\newlist{parts}{enumerate}{2}
\setlist[parts,1]{label=(\alph*),leftmargin=*,itemsep=5pt,topsep=3pt}
\setlist[parts,2]{label=(\roman*),leftmargin=*,itemsep=3pt,topsep=2pt}

\newcommand{\pts}[1]{\hfill   \textit{[#1]}}


\begin{document}

\begin{center}
  {\large\bfseries BANARAS HINDU UNIVERSITY}\\[2pt]
  {\normalsize B.Sc. (Hons.) --- Semester II Examination, 2013-14 (Ancillary)}\\[6pt]
  \rule{0.8\linewidth}{0.5pt}\\[6pt]
  {\large\bfseries Paper No. BSC-06A: Ancillary Mathematics (Old Course)}\\[4pt]
  {\normalsize Time: 3 Hours \hfill Full Marks: 70}
\end{center}
\rule{\linewidth}{0.4pt}
\smallskip



\noindent   \textit{Instructions: Answer   \textbf{five} questions including Question~1 which is compulsory. Figures in right-hand margin indicate marks.}
\bigskip




\noindent   \textbf{Question 1.}   \textit{(Compulsory --- 2 marks each.)}

\begin{parts}
  \item Draw the graph of the function $y = e^x$.
  \item Find $\lim_{x  o 1} \frac{x-1}{x^2-1}$.
  \item Find $\frac{d}{dx}(\log \sqrt{x})$.
  \item Find $\int (ax^2 + 4x^3) \, dx$.
  \item Define the order and the degree of a differential equation.
  \item If $A = 
\begin{pmatrix} 1 & 2 & 3 \\ 2 & 4 & 1 \\ 3 & 0 & 5 \end{pmatrix}$, then find the transpose of $A$.
  \item If $A = 
\begin{pmatrix} 1 & 1 \\ 3 & 0 \\ 3 & 1 \end{pmatrix}$ and $B = 
\begin{pmatrix} 1 & -1 \\ 2 & 2 \\ -2 & 0 \end{pmatrix}$, find $2A - B$.
  \item Let $A = 
\begin{pmatrix} 1 & 3 \\ 2 & 4 \end{pmatrix}$ and $B = 
\begin{pmatrix} 2 & 0 \\ 1 & 3 \end{pmatrix}$. Find $AB$.
  \item Define lower and upper triangular matrices.
  \item Solve: $\frac{dy}{dx} = \frac{x+y}{x}$.
  \item Define diagonal and scalar matrices.
\end{parts}

\medskip\rule{\linewidth}{0.2pt}

\medskip


\noindent   \textbf{Question 2.} Find the following: \pts{12}

\begin{parts}
  \item $\frac{d}{dx}(2^x)$
  \item $\frac{d}{dx}(x^n \cos x)$
  \item $\frac{d}{dx}(\sin \sqrt{x})$
  \item $\frac{d}{dx}(x \log x)$
\end{parts}

\medskip\rule{\linewidth}{0.2pt}

\medskip


\noindent   \textbf{Question 3.} Find the following: \pts{12}

\begin{parts}
  \item $\int (1-3x)(1+x) \, dx$
  \item $\int \frac{3}{x} \, dx$
  \item $\int_0^1 \frac{1}{x^2+1} \, dx$
  \item $\int x e^x \, dx$
\end{parts}

\medskip\rule{\linewidth}{0.2pt}

\medskip


\noindent   \textbf{Question 4.}

\begin{parts}
  \item Find the turning points on the curve $y = x^3 - 3x^2 + 3x$ and distinguish them. Also find the maximum and minimum values of the function. \pts{6}
  \item Show that $y = \frac{c}{x}$ is a solution of the differential equation $x \frac{d^2y}{dx^2} + \frac{dy}{dx} = 0$. \pts{6}
\end{parts}

\medskip\rule{\linewidth}{0.2pt}

\medskip


\noindent   \textbf{Question 5.}

\begin{parts}
  \item Find the area between the curve $y = x(4-x)$ and the $x$-axis from $x=0$ to $x=5$. \pts{6}
  \item Show that $
\begin{vmatrix} 1+a & b & c \\ a & 1+b & c \\ a & b & 1+c \end{vmatrix} = 1 + a + b + c$. \pts{6}
\end{parts}

\medskip\rule{\linewidth}{0.2pt}

\medskip


\noindent   \textbf{Question 6.}

\begin{parts}
  \item Compute the inverse of the following matrix:
  \[
  
\begin{pmatrix}
  1 & -3 & 2 \\
  2 & 0 & 0 \\
  1 & 4 & 1
  \end{pmatrix}
  \] \pts{6}
  \item If $A = 
\begin{pmatrix} 3 & -4 \\ 2 & 1 \end{pmatrix}$ and $B = 
\begin{pmatrix} 2 & 1 \\ 1 & 3 \end{pmatrix}$, show that $(AB)' = B' A'$. \pts{6}
\end{parts}

\medskip\rule{\linewidth}{0.2pt}

\medskip


\noindent   \textbf{Question 7.}

\begin{parts}
  \item Solve the following system of equations by matrix method:
  
\begin{align*}
  2x_1 - x_2 + 3x_3 &= 9 \\
  x_1 + x_2 + x_3 &= 6 \\
  x_1 - x_2 + x_3 &= 2
  \end{align*} \pts{6}
  \item Solve the following system of equations by Cramer's rule:
  
\begin{align*}
  3x + 5y - 7z &= 13 \\
  4x + y - 12z &= 6 \\
  2x + 9y - 3z &= 20
  \end{align*} \pts{6}
\end{parts}

\vfill
\rule{\linewidth}{0.4pt}

\begin{center}\small
  \href{https://bhu-exam.vercel.app/}{Click here to visit our website}\\[4pt]
  \href{https://me-aryan.vercel.app/}{Click here to visit our contributor}\\[4pt]
  \href{https://quashan-me.vercel.app/}{Click here to visit our contributor}
\end{center}

\end{document}
